\documentclass[11pt]{article}
\usepackage[a4paper,margin=28mm]{geometry}
\usepackage{amsmath,amssymb,amsthm}
\usepackage[colorlinks=true,linkcolor=blue,urlcolor=blue,citecolor=blue]{hyperref}
\newtheorem*{theorem*}{Theorem}
\newtheorem*{corollary*}{Corollary}
\newenvironment{statement}[1]{\par\medskip\noindent\phantomsection\textbf{#1}\ }{\par\medskip}
\setlength{\emergencystretch}{3em}
\newcommand{\Q}{\mathbb Q}
\newcommand{\Z}{\mathbb Z}
\newcommand{\F}{\mathbb F}
\DeclareMathOperator{\End}{End}
\DeclareMathOperator{\Gal}{Gal}
\DeclareMathOperator{\GL}{GL}
\DeclareMathOperator{\PGL}{PGL}
\DeclareMathOperator{\Hom}{Hom}
\DeclareMathOperator{\Res}{Res}
\DeclareMathOperator{\Ind}{Ind}
\DeclareMathOperator{\tr}{tr}
\DeclareMathOperator{\disc}{disc}
\title{No rational point of order \(45\) on a geometrically split abelian surface simple over \(\Q\)}
\author{}
\date{14 November 2026}

\begin{document}
\maketitle

\begin{theorem*}
Let \(A/\Q\) be a principally polarized abelian surface that is simple over \(\Q\) and
geometrically split. Then \(A(\Q)\) has no point of order \(45\).
\end{theorem*}

\begin{corollary*}
No genus-\(2\) Jacobian over \(\Q\) that is geometrically split and simple over \(\Q\) has a
rational torsion subgroup containing an element of order \(45\). In particular
\href{https://genus-2-torsion.github.io/group.html?g=45}{\([45]\)} does not occur as the rational
torsion subgroup of such a Jacobian.
\end{corollary*}

Here ``geometrically split'' means that \(A_{\bar\Q}\) is isogenous to a product of two elliptic
curves. We impose no restriction on the degree of this isogeny, on the field of definition of the
geometric endomorphisms or on the integral endomorphism ring. The principal polarization and the
simplicity over \(\Q\) are used only in Section~\ref{parents}. The theorem says nothing about
surfaces split over \(\Q\): a product of elliptic curves over \(\Q\) with rational points of orders
\(5\) and \(9\) has a rational point of order \(45\), and the census contains Jacobians split over
\(\Q\) with torsion subgroup \([45]\). For geometric squares we prove a stronger statement, which
needs neither simplicity over \(\Q\) nor a polarization:

\begin{statement}{Proposition A.}\label{prop-A}
Let \(A/\Q\) be an abelian surface such that \(A_{\bar\Q}\) is isogenous to the square of an
elliptic curve. Then \(A(\Q)\) has no point of order \(45\).
\end{statement}

\emph{Outline.} Section~\ref{lemmas} proves two lemmas on elliptic curves: coprime torsion can be
combined in an isogeny class, and when \(E[p]^{\mathrm{ss}}\) is a sum of two distinct characters,
every Galois-stable lattice in \(V_p(E)^r\) splits into lattices of rank \(2\).
Section~\ref{square} proves Proposition~A: CM squares by a point count over \(\F_2\), non-CM
squares by the residual-representation lemmas of \cite{P63} at the primes \(5\) and \(3\) and, when
the endomorphism field is quadratic, by moving rational \(9\)-torsion to an elliptic curve over that
field. In the remaining case \(A_{\bar\Q}\) is isogenous to \(E\times F\) with \(E\) and \(F\) not
isogenous. Then \(E\) and \(F\) are conjugate over a quadratic field \(K\), and the principal
polarization identifies each of them with the quotient of \(A_K\) by the other. Section~\ref{parents}
shows that \(E(K)_{\mathrm{tors}}\cong\Z/15\), that \(A_K\) is an \((n,n)\)-gluing of \(E\) and
\(E^\sigma\) with \(v_3(n)=1\), and that its \(3\)-part is a conjugate \((3,3)\)-gluing \(B\) over
\(\Q\) with a rational point of order \(9\). Section~\ref{conditions} turns this into two conditions
on \(E\): the norm of its discriminant is a rational cube, and its flex coefficient \(a_3\)
satisfies \(a_3/\sigma(a_3)\in K^{\times3}\). Section~\ref{families} shows that no quadratic point
of \(X_1(15)\) satisfies both. On each of the four families of quadratic points, the second
condition forces a primitive value of a cyclotomic binary quartic to be a cube or five times a cube;
a descent in \(\Q(\zeta_5)\) following Poonen and elliptic Chabauty over \(\Q(\zeta_5)\)
(Section~\ref{cyclotomic}) leave four parameter values, where the point is a cusp or the first
condition fails.

\section{Conventions and inputs}\label{inputs}

All number fields are subfields of a fixed algebraic closure \(\bar\Q\). For a number field \(k\)
we write \(G_k=\Gal(\bar\Q/k)\), and \(\bar\chi_\ell\) for the mod-\(\ell\) cyclotomic character.
For an abelian variety \(X\), \(T_\ell(X)\) is the Tate module, \(V_\ell(X)=T_\ell(X)\otimes\Q_\ell\),
\(X[n]\) is the kernel of multiplication by \(n\) and \(\End^0=\End\otimes\Q\). For an elliptic curve
\(E\), \(e_n\) is the Weil pairing on \(E[n]\); an \emph{anti-isometry} \(E[n]\to E'[n]\) is an
isomorphism \(\psi\) with \(e_n(\psi x,\psi y)=e_n(x,y)^{-1}\). For a quadratic field \(K\) we write
\(\sigma\) for its nontrivial automorphism and \(\tilde\sigma\in G_{\Q}\) for a lift of \(\sigma\);
for a variety \(Y\) over \(K\), \(Y^\sigma\) is the conjugate variety and
\(\sigma:Y(\bar\Q)\to Y^\sigma(\bar\Q)\) the bijection induced by \(\tilde\sigma\)
\cite[Section~1]{P2214}. We write \(k^{\times3}\) for the nonzero cubes of a field \(k\), and
\(v_\ell\) for the \(\ell\)-adic valuation. We use the functoriality of the Weil pairing and Galois
descent for morphisms of varieties without further comment.

\begin{itemize}
\item \textbf{(I1) Torsion of elliptic curves over \(\Q\) and over quadratic fields.} By Mazur
\cite[Theorem~(8)]{Ma}, no elliptic curve over \(\Q\) has a rational point of order \(15\). By
Kenku--Momose \cite[Theorem~(0.1)]{KM} and Kamienny \cite{Ka} (see the statement in
\cite[Section~1]{Na}), the torsion subgroup of an elliptic curve over a quadratic field is
isomorphic to \(\Z/m\) with \(1\le m\le16\) or \(m=18\), to \(\Z/2\times\Z/2m\) with \(1\le m\le6\),
to \(\Z/3\times\Z/3m\) with \(m\in\{1,2\}\), or to \(\Z/4\times\Z/4\). In particular there is no
point of order \(45\), and a torsion subgroup containing a point of order \(15\) is isomorphic to
\(\Z/15\). Over \(\Q(\sqrt{-3})\) only Mazur's groups, \(\Z/3\times\Z/3\) and \(\Z/3\times\Z/6\)
occur \cite[Theorem~2(ii)]{Na}; in particular there is no point of order \(15\). (The finite checks
in this item are part~A of the certificate.)

\item \textbf{(I2) Reduction} \cite{ST}, as stated in \cite[Section~1]{P63}: let \(K'\) be a finite
extension of \(\Q_p\) with residue field \(k\) and \(X\) an abelian variety over \(K'\). (a) Good
reduction is detected by the inertia action on \(T_\ell(X)\), \(\ell\neq p\)
\cite[\S1, Theorem~1]{ST}; in particular it is preserved by isogenies \cite[\S1, Corollary~2]{ST}.
(b) If \(X\) has good reduction \(\widetilde X\) and \(m\) is prime to \(p\), then \(X(K')[m]\)
injects into \(\widetilde X(k)\) \cite[Lemma~2]{ST}. (c) If \(X\) has potentially good reduction,
then it has good reduction over a finite extension of \(K'\) with residue field \(k\)
\cite[\S2]{ST}. (d) (Weil.) An abelian surface \(Y\) over \(\F_q\) has at most \((1+\sqrt q)^4\)
points.

\item \textbf{(I3) The case distinction} \cite[Sections~2 and~5]{P63}. A geometrically split abelian
surface \(A/\Q\) is of exactly one of the types \textbf{(N)} (\(A_{\bar\Q}\) is isogenous to
\(E_1\times E_2\) with \(E_1\), \(E_2\) not isogenous), \textbf{(CM)} (\(A_{\bar\Q}\) is isogenous to
\(E^2\) with \(E\) a CM curve) and \textbf{(S)} (\(A_{\bar\Q}\) is isogenous to \(E^2\) with \(E\)
without CM). In case (S), \(D=\End^0(A_{\bar\Q})\cong M_2(\Q)\); \(L\) is the smallest field over
which all geometric endomorphisms are defined; \(\Gamma=\Gal(L/\Q)\) embeds in \(\PGL_2(\Q)\), the
element \(g\) acting on \(D\) by conjugation with \(R_g\in D^\times\); and \(\Gamma\) is one of
\(C_1\), \(C_2\), \(C_3\), \(C_4\), \(C_6\), \(V_4\), \(D_6\), \(D_8\), \(D_{12}\) (\(D_{2m}\) is
dihedral of order \(2m\)). An element of \(\Gamma\) of order \(m\in\{3,4,6\}\) has a unique
representative \(R\) with \(\tr R=1\), and \(R^2-R+1/c=0\) with \(c=1,2,3\) respectively.

\item \textbf{(I4)} \cite[Lemma~2]{P63}. Let \(X\) be an abelian surface over a field \(k\),
\(C\subset X\) an elliptic curve defined over \(k\), and \(Q\in X(k)\) of prime order \(q\). Then
either \(Q\in C(k)\), or the image of \(Q\) in \((X/C)(k)\) has order \(q\).

\item \textbf{(I5)} \cite[Lemma~3]{P63}. In case (S), every elliptic curve \(C\subset A_{\bar\Q}\) is
defined over \(L\), any two of them are \(L\)-isogenous, and \(A_L/C\) is \(L\)-isogenous to \(C\).

\item \textbf{(I6) Residual representations} \cite[Section~6]{P63}. Let \(A\) be of type (S), \(p\)
an odd prime with \(A(\Q)[p]\neq0\), and \(D_p=D\otimes\Q_p\).
\begin{enumerate}
\item[(a)] (\emph{Unramified order}, \cite[Lemma~6]{P63}.) Suppose that a \(\Z_p\)-order
\(\mathcal B\subset D_p\), isomorphic to the ring of integers of the unramified quadratic extension
of \(\Q_p\), is preserved by \(\Gamma\), and that some \(g_0\in\Gamma\) has a representative
\(R_{g_0}\in\mathcal B^\times\) whose reduction in \(\mathcal B/p\mathcal B\cong\F_{p^2}\) is not in
\(\F_p\). Then \(\Q(\zeta_p)\subseteq L\).
\item[(b)] (\emph{Split blocks}, \cite[Lemma~7]{P63}.) Suppose that
\(\mathcal B=\Z_pe_1\oplus\Z_pe_2\subset D_p\), with orthogonal idempotents \(e_1,e_2\) of rank one
and \(e_1+e_2=1\), is preserved by \(\Gamma\), and that some element of \(\Gamma\) exchanges \(e_1\)
and \(e_2\). Let \(\Gamma_0\) be the stabilizer of \(e_1\) and \(K_0=L^{\Gamma_0}\). If the character
\(\eta\) by which \(G_{K_0}\) acts on the line \(e_1D_pe_2\) has nontrivial reduction modulo \(p\),
then \(\Q(\zeta_p)\subseteq L\).
\end{enumerate}
In (a) the proof replaces \(T_p(A)\) by the lattice generated over the \(\Gamma\)-stable maximal
order \(\End_{\Z_p}(\Lambda)\supset\mathcal B\), where \(\Lambda\subset\Q_p^2\) is a
\(\mathcal B\)-stable lattice, and in (b) by \(\mathcal B\,T_p(A)\). Both proofs end with the
polarization duality of \cite[Lemma~5]{P63} (\cite[Lemma~3]{P30}), which needs only a polarization of
\(A\) over \(\Q\), of any degree. The proofs use nothing else about \(A\): neither the simplicity
over \(\Q\) nor the point of order \(63\) assumed in \cite[Section~2]{P63}. The same arguments appear
in \cite[Lemma~4 and Section~8]{P30}.

\item \textbf{(I7) Weil restriction} \cite[\S1]{Mi}. Let \(K\) be a quadratic field and \(E\) an
elliptic curve over \(K\). The Weil restriction \(\Res_{K/\Q}E\) is an abelian surface over \(\Q\)
with \((\Res_{K/\Q}E)(\Q)=E(K)\) and \((\Res_{K/\Q}E)_K\cong E\times E^\sigma\). Its
\(\bar\Q\)-points are the pairs \((x,y)\) with \(x\in E(\bar\Q)\) and \(y\in E^\sigma(\bar\Q)\); an
element \(g\in G_K\) maps \((x,y)\) to \((gx,gy)\), an element \(g\notin G_K\) maps it to
\((gy,gx)\), and the rational points are the pairs \((P,\sigma P)\) with \(P\in E(K)\).

\item \textbf{(I8) Complements and gluing} \cite[Section~5.3 and Corollary~12.1.5]{BL}, \cite{FK},
\cite{Ka97}. For an elliptic curve \(C\) in an abelian surface \(X\) with a polarization \(\lambda\),
the complement of \(C\) is the connected component of the kernel of \(i_C^\vee\circ\lambda\), where
\(i_C\) is the inclusion. A principally polarized abelian surface containing two complementary
elliptic curves \(E\), \(F\) is the quotient of \(E\times F\) by the graph of an anti-isometry
\(E[n]\to F[n]\) (an \((n,n)\)-gluing). We prove the parts of this that we use in
Section~\ref{parents}.

\item \textbf{(I9) Tate normal form} \cite[Table~3, the entry for \(\Z/5\)]{Ku}. Let \(k\) be a field
of characteristic \(0\) and put \(E_d:\ y^2+(1-d)xy-dy=x^3-dx^2\), whose discriminant is
\(d^5(d^2-11d-1)\). Every pair \((E,P)\) of an elliptic curve over \(k\) and a point \(P\in E(k)\) of
order \(5\) is isomorphic over \(k\) to \((E_d,(0,0))\) for some \(d\in k\) with
\(d^5(d^2-11d-1)\neq0\). Kubert works over \(\Q\), but his derivation from the Tate normal form uses
only field operations.
\end{itemize}

\section{Two lemmas on elliptic curves}\label{lemmas}

\begin{statement}{Lemma 1 (combining coprime torsion).}\label{lemma-coprime}
Let \(k\) be a number field and \(m\), \(n\) coprime positive integers. If a \(k\)-isogeny class of
elliptic curves over \(k\) contains a curve with a \(k\)-rational point of order \(m\) and a curve
with a \(k\)-rational point of order \(n\), then it contains a curve with a \(k\)-rational point of
order \(mn\).
\end{statement}

\emph{Proof.} This is \cite[Lemma~1]{P63}. Let \(f:E_m\to E_n\) be a \(k\)-isogeny between the two
curves and \(H\) the subgroup of \(\ker f\) of elements whose order involves only primes dividing
\(n\); it is characteristic, hence defined over \(k\), and \(f=b\circ a\) with
\(a:E_m\to E'=E_m/H\). As \(\deg a\) is prime to \(m\), \(a\) maps the point of order \(m\) to a point
of order \(m\) in \(E'(k)\); as \(\deg b\) is prime to \(n\), the dual isogeny \(\hat b\) maps the
point of order \(n\) to a point of order \(n\) in \(E'(k)\). Their sum has order \(mn\). \qed

\begin{statement}{Lemma 2 (stable lattices).}\label{lemma-lattice}
Let \(k\) be a number field, \(p\) a prime, and \(E/k\) an elliptic curve such that
\(E[p]^{\mathrm{ss}}\cong\alpha\oplus\beta\) for two distinct characters
\(\alpha,\beta:G_k\to\F_p^\times\). Put \(T=T_p(E)\) and \(V=V_p(E)\), and let \(r\ge1\).
\begin{enumerate}
\item[(a)] Every \(G_k\)-stable \(\Z_p\)-lattice in \(V^r\) is a direct sum of \(r\) \(G_k\)-stable
lattices of rank \(2\), each contained in a \(G_k\)-subrepresentation of \(V^r\) isomorphic to \(V\).
\item[(b)] If an abelian variety \(Y\) over \(k\) is \(k\)-isogenous to \(E^r\), then there are
elliptic curves \(E_1,\dots,E_r\) over \(k\), each \(k\)-isogenous to \(E\), such that
\(Y[p^\infty]\cong\bigoplus_jE_j[p^\infty]\) as \(G_k\)-modules. In particular,
if \(Y(k)\) has a point of order \(p^e\), then so has some \(E_j(k)\).
\end{enumerate}
\end{statement}

\emph{Proof.} Let \(\mathcal R\) be the \(\Z_p\)-span of the image of \(G_k\) in
\(\End_{\Z_p}(T)\cong M_2(\Z_p)\). It is a finitely generated, hence free, \(\Z_p\)-submodule of
\(\End_{\Z_p}(T)\), and a subring containing \(1\). Choose \(g\in G_k\) with
\(\alpha(g)\neq\beta(g)\). The characteristic polynomial of \(g\) on \(T\) has the distinct roots
\(\alpha(g)\), \(\beta(g)\) modulo \(p\), so by Hensel's lemma it has roots \(\lambda_1,\lambda_2\in\Z_p\)
with \(\lambda_1-\lambda_2\in\Z_p^\times\). By the Cayley--Hamilton theorem,
\(\epsilon_1=(g-\lambda_2)/(\lambda_1-\lambda_2)\) and \(\epsilon_2=1-\epsilon_1\) are orthogonal
idempotents, and they lie in \(\mathcal R\). They have rank one, since \(g\) is not a scalar, so
\(T=\Z_pv_1\oplus\Z_pv_2\) with \(\epsilon_iT=\Z_pv_i\). In this basis \(\epsilon_i\mathcal R\epsilon_i\)
contains \(\epsilon_i\) and is contained in \(\Z_p\epsilon_i\), and the corners
\(\epsilon_1\mathcal R\epsilon_2\), \(\epsilon_2\mathcal R\epsilon_1\) are \(\Z_p\)-submodules of the
off-diagonal lines. As \(\epsilon_1+\epsilon_2=1\), \(\mathcal R\) is the sum of its four corners, so
\[\mathcal R=\begin{pmatrix}\Z_p&p^a\Z_p\\ p^b\Z_p&\Z_p\end{pmatrix}\qquad\text{with }
a,b\in\{0,1,2,\dots\}\cup\{\infty\},\ p^\infty=0.\]

(a) Write \(V^r=V\otimes_{\Q_p}U\) with \(U=\Q_p^r\), where \(G_k\) acts on the first factor. A
\(G_k\)-stable lattice \(T'\subset V^r\) is stable under \(\mathcal R\otimes1\), so
\(T'=\epsilon_1T'\oplus\epsilon_2T'\) with \(\epsilon_iT'=v_i\otimes\Lambda_i\) for lattices
\(\Lambda_1,\Lambda_2\subset U\). The off-diagonal elements of \(\mathcal R\) give
\(p^a\Lambda_2\subseteq\Lambda_1\) and \(p^b\Lambda_1\subseteq\Lambda_2\), and conversely every lattice
\(v_1\otimes\Lambda_1\oplus v_2\otimes\Lambda_2\) with these two inclusions is \(\mathcal R\)-stable,
hence \(G_k\)-stable. By the elementary divisor theorem there are a basis \(w_1,\dots,w_r\) of
\(\Lambda_1\) and integers \(d_1,\dots,d_r\) with \(\Lambda_2=\bigoplus_j p^{d_j}\Z_pw_j\). The two
inclusions say \(a+d_j\ge0\) and \(d_j\le b\) for all \(j\), so each
\(T'_j=\Z_p(v_1\otimes w_j)\oplus p^{d_j}\Z_p(v_2\otimes w_j)\) is \(G_k\)-stable, and
\(T'=\bigoplus_jT'_j\) with \(T'_j\subset V\otimes w_j\cong V\).

(b) An isogeny identifies \(V_p(Y)\) with \(V^r\) \(G_k\)-equivariantly, and \(T_p(Y)\) with a lattice
\(T'=\bigoplus_jT'_j\) as in (a). Identify \(V\otimes w_j\) with \(V\) and choose \(c_j\ge0\) with
\(p^{c_j}T'_j\supseteq T\). Then \(C_j=p^{c_j}T'_j/T\) is a finite \(G_k\)-stable subgroup of
\(V/T=E[p^\infty]\), and \(E_j=E/C_j\) is an elliptic curve over \(k\), \(k\)-isogenous to \(E\),
with \(T_p(E_j)=p^{c_j}T'_j\cong T'_j\). Hence
\(Y[p^\infty]=V_p(Y)/T_p(Y)\cong\bigoplus_j(V\otimes w_j)/T'_j\cong\bigoplus_jE_j[p^\infty]\). Taking
\(G_k\)-invariants, \(Y(k)[p^\infty]\cong\bigoplus_jE_j(k)[p^\infty]\), and a point of order \(p^e\)
in the sum has a coordinate of order \(p^e\). \qed

\medskip
The lemma needs neither the irreducibility of \(V\) nor any hypothesis on complex multiplication:
only the two idempotents \(\epsilon_1,\epsilon_2\) of \(\mathcal R\) are used. They exist because
\(\alpha\neq\beta\); this is why the field \(\Q(\sqrt{-3})\) has to be excluded in the case
\(\Gamma=C_2\) of Section~\ref{square}.

\section{The geometric-square case}\label{square}

\emph{Proof of \hyperref[prop-A]{Proposition A}.} Let \(A_{\bar\Q}\) be isogenous to \(E^2\), and
suppose that \(R_{45}\in A(\Q)\) has order \(45\). Put \(R_9=5R_{45}\), \(R_5=9R_{45}\) and
\(R_3=15R_{45}\), of orders \(9\), \(5\) and \(3\).

\emph{The CM case.} Suppose that \(E\) has complex multiplication. As in
\cite[Section~4]{P63}, \(E\) has potentially good reduction at \(2\), because its \(j\)-invariant is an
algebraic integer \cite[\S2]{ST}; so \(A_{\Q_2}\) acquires good reduction over a finite extension by
(I2)(a), and by (I2)(c) over a finite extension \(K'\) of \(\Q_2\) with residue field \(\F_2\). By
(I2)(b) the subgroup generated by \(R_{45}\), of odd order, injects into the reduction, whose number
of points is at most \((1+\sqrt2)^4=17+12\sqrt2<34\) by (I2)(d). This contradicts \(45\le34\).
(\cite[Proposition~2]{P30} shows more: \(A(\Q_2)\) has no point of order \(15\).)

\emph{The non-CM case.} Now \(A\) is of type (S), with \(D\), \(L\), \(\Gamma\) as in (I3). Since
\(R_5\) and \(R_3\) are rational, (I6) applies with \(p=5\) and with \(p=3\). Modulo \(5\) the
discriminants \(-3\), \(-1\), \(-\tfrac13\) of the polynomials \(X^2-X+\tfrac1c\) of (I3) are \(2\),
\(4=2^2\) and \(3\); the nonzero squares modulo \(5\) are \(1\) and \(4\). The field \(\Q(\zeta_5)\) is
cyclic of degree \(4\), so if \(\Q(\zeta_5)\subseteq L\), then \(\Gamma\) has a quotient isomorphic
to \(C_4\). We go through the groups of (I3).
\begin{enumerate}
\item[(i)] \(\Gamma=C_1\). Then \(L=\Q\), and by (I5) every elliptic curve \(C\subset A\) is defined
over \(\Q\), with \(A/C\) \(\Q\)-isogenous to \(C\). By (I4), applied to \(R_3\) and \(R_5\), the
\(\Q\)-isogeny class of \(C\) contains curves with rational points of orders \(3\) and \(5\), and
\hyperref[lemma-coprime]{Lemma~1} gives an elliptic curve over \(\Q\) with a rational point of order
\(15\), contrary to (I1).

\item[(ii)] \(\Gamma\in\{C_3,C_6,D_6,D_{12}\}\). Let \(r\) be a rotation of order \(3\) or \(6\)
generating the rotation subgroup and \(R\) its trace-one representative, with \(c\in\{1,3\}\). Every
element of \(\Gamma\) maps \(R\) to \(R\) or \(1-R\) \cite[Section~7]{P63}, so
\(\mathcal B=\Z_5[R]\) is preserved by \(\Gamma\). Its discriminant is a non-square unit modulo \(5\),
so \(\mathcal B\) is the ring of integers of the unramified quadratic extension of \(\Q_5\);
\(R\in\mathcal B^\times\) because \(\det R=1/c\) is a unit, and its reduction is a root of an
irreducible quadratic over \(\F_5\). By (I6)(a), \(\Q(\zeta_5)\subseteq L\). But \(C_3\), \(C_6\) and
the abelianizations \(C_2\) and \(C_2\times C_2\) of \(D_6\) and \(D_{12}\) have no quotient \(C_4\).

\item[(iii)] \(\Gamma=D_8\). Let \(r\) be a rotation of order \(4\) and \(R\) its trace-one
representative, \(R^2-R+\tfrac12=0\). Again \(\mathcal B=\Z_5[R]\) is preserved by \(\Gamma\). The
discriminant \(-1\equiv2^2\) is a square modulo \(5\), so the roots \(\lambda_1,\lambda_2\) of
\(X^2-X+\tfrac12\) lie in \(\Z_5^\times\), with \(\lambda_1-\lambda_2\) a unit, and
\(\mathcal B=\Z_5e_1\oplus\Z_5e_2\) with \(e_1=(R-\lambda_2)/(\lambda_1-\lambda_2)\) and
\(e_2=1-e_1\). The rotations fix \(e_1\) and \(e_2\), and the reflections, which map \(R\) to
\(1-R\), exchange them because \(\lambda_1+\lambda_2=1\). The stabilizer of \(e_1\) is the rotation
subgroup, and \(RXR^{-1}=(\lambda_1/\lambda_2)X\) for \(X\in e_1D_5e_2\), so \(\eta\) maps a lift of
\(r\) to \(\lambda_1/\lambda_2\), a primitive fourth root of unity, whose reduction modulo \(5\) still
has order \(4\). By (I6)(b), \(\Q(\zeta_5)\subseteq L\); but the abelianization of \(D_8\) is
\(C_2\times C_2\).

\item[(iv)] \(\Gamma=V_4\). The map \([M]\mapsto v_5(\det M)\bmod2\) is a homomorphism
\(\Gamma\to\Z/2\); choose \(r\neq1\) in its kernel and a representative \(R\) of trace \(0\),
scaled by a power of \(5\) so that \(\det R\in\Z_5^\times\). Then \(R^2=a\) with
\(a\in\Q^\times\cap\Z_5^\times\). As \(\Gamma\) is abelian, every \(R_g\) conjugates \(R\) to
\(\pm R\), so \(\mathcal B=\Z_5[R]\) is preserved by \(\Gamma\). If \(a\) is not a square modulo
\(5\), then \(\mathcal B\) is the unramified quadratic order, \(R\in\mathcal B^\times\), its reduction
is not in \(\F_5\), and (I6)(a) applies. If \(a=b^2\) with \(b\in\Z_5^\times\), then
\(\mathcal B=\Z_5e_+\oplus\Z_5e_-\) with \(e_\pm=(1\pm R/b)/2\). Some element of \(\Gamma\) conjugates
\(R\) to \(-R\): otherwise \(\Gamma\) would embed in the image of the centralizer
\(\Q_5[R]^\times\cong\Q_5^\times\times\Q_5^\times\) in \(\PGL_2(\Q_5)\), which is
\(\cong\Q_5^\times\) and has only cyclic finite subgroups. That element exchanges \(e_+\) and \(e_-\);
the stabilizer of \(e_+\) is \(\{1,r\}\), and \(RXR^{-1}=-X\) for \(X\in e_+D_5e_-\), so \(\eta\) is the
nontrivial quadratic character and (I6)(b) applies. In both cases \(\Q(\zeta_5)\subseteq L\), which
is impossible because \(L\) is biquadratic.

\item[(v)] \(\Gamma=C_4\). Here we use the prime \(3\). Let \(r\) generate \(\Gamma\) and let \(R\)
be its trace-one representative, \(R^2-R+\tfrac12=0\). The discriminant \(-1\equiv2\) is not a square
modulo \(3\), so \(\mathcal B=\Z_3[R]\) is the ring of integers of the unramified quadratic extension
of \(\Q_3\); it is fixed by \(\Gamma=\langle r\rangle\); \(R\in\mathcal B^\times\) since
\(\det R=\tfrac12\) is a \(3\)-adic unit; and its reduction is a root of the irreducible polynomial
\(X^2-X+2\) over \(\F_3\). By (I6)(a), \(\Q(\zeta_3)=\Q(\sqrt{-3})\subseteq L\). But a cyclic quartic
field \(L\) does not contain \(\Q(\sqrt{-3})\): complex conjugation would restrict to an element of
\(\Gal(L/\Q)\cong C_4\) that acts nontrivially on the unique quadratic subfield, hence to an element
of order \(4\), whereas complex conjugation has order at most \(2\). (Equivalently, \(L\) is totally
real by the argument of \cite[Section~9.1]{P63}, which uses only that every element of \(\Gamma\) is
the class of a power \(R^k\) with \(\det R^k=2^{-k}>0\).)

\item[(vi)] \(\Gamma=C_2\). Then \(L=K\) is a quadratic field. By (I5), an elliptic curve
\(C\subset A_{\bar\Q}\) is defined over \(K\), and \(A_K/C\) is \(K\)-isogenous to \(C\); with a
complement of \(C\), also \(K\)-isogenous to \(C\), we see that \(A_K\) is \(K\)-isogenous to
\(C^2\). By (I4), applied over \(K\) to \(R_3\) and \(R_5\), and by
\hyperref[lemma-coprime]{Lemma~1}, the \(K\)-isogeny class of \(C\) contains a curve \(E_0\) with a
\(K\)-rational point of order \(15\). By (I1), \(E_0(K)_{\mathrm{tors}}\cong\Z/15\) and
\(K\neq\Q(\sqrt{-3})\). The semisimplification of \(E_0[3]\) is \(\mathbf1\oplus\bar\chi_3\) (a
rational point of order \(3\) and the Weil pairing), and \(\bar\chi_3\) is nontrivial on \(G_K\)
because \(K\neq\Q(\sqrt{-3})\). As \(A_K\) is \(K\)-isogenous to \(E_0^2\),
\hyperref[lemma-lattice]{Lemma~2}(b) with \(p=3\) and \(r=2\), applied to the point \(R_9\in A(K)\) of
order \(9\), gives an elliptic curve \(E_2\), \(K\)-isogenous to \(E_0\), with a \(K\)-rational point
of order \(9\). \hyperref[lemma-coprime]{Lemma~1}, with \(E_2\) and \(E_0\), gives an elliptic curve
over \(K\) with a \(K\)-rational point of order \(45\), contrary to (I1).
\end{enumerate}
This covers all groups in (I3). \qed

\medskip
No step uses simplicity over \(\Q\), and the only polarization used is a polarization over \(\Q\) of
arbitrary degree, inside (I6).

\section{Nonisogenous factors: the elliptic curves and the gluing}\label{parents}

\begin{statement}{Lemma 3.}\label{lemma-parents}
Let \(A/\Q\) be a principally polarized abelian surface that is simple over \(\Q\) and of type (N),
and suppose that \(A(\Q)\) has a point of order \(45\). Then there are a quadratic field \(K\), an
elliptic curve \(E\) over \(K\) and a point \(P\in E(K)\) of order \(3\) with the following
properties.
\begin{enumerate}
\item[(a)] \(E(K)_{\mathrm{tors}}\cong E^\sigma(K)_{\mathrm{tors}}\cong\Z/15\).
\item[(b)] There is a \(G_K\)-equivariant anti-isometry \(\psi:E[3]\to E^\sigma[3]\) with
\(\psi(P)=-\sigma(P)\).
\item[(c)] There are an abelian surface \(B\) over \(\Q\), isogenies
\(\varphi:\Res_{K/\Q}E\to B\) and \(\Phi:B\to\Res_{K/\Q}E\) over \(\Q\) with \(\Phi\circ\varphi=[3]\),
such that \(\ker\varphi\) is the graph of \(\psi\) (over \(K\)), and a point \(R\in B(\Q)\) of order
\(9\) with \(\Phi(R)=(P,\sigma P)\).
\end{enumerate}
\end{statement}

\emph{Proof.} Let \(\lambda:A\to A^\vee\) be the principal polarization, which is defined over
\(\Q\), and let \(R_{45}\in A(\Q)\) have order \(45\).

\emph{The two elliptic curves.} As in \cite[Section~3]{P63}, \(A_{\bar\Q}\) contains exactly two
elliptic curves, and \(G_{\Q}\) permutes them. If it fixed them, they would be defined over \(\Q\) and
\(A\) would not be simple over \(\Q\). So they are defined over a quadratic field \(K\) and exchanged
by \(\tilde\sigma\); call them \(E\) and \(F=\tilde\sigma(E)=E^\sigma\).

\emph{Each is the quotient by the other.} For an elliptic curve \(C\subset A_K\) let
\(i_C:C\to A_K\) be the inclusion and \(\lambda^0_C:C\to C^\vee\) the canonical principal
polarization. The dual \(i_E^\vee:A_K^\vee\to E^\vee\) of the inclusion is surjective with connected
kernel \((A_K/E)^\vee\), so, \(\lambda\) being an isomorphism, the kernel of \(i_E^\vee\circ\lambda\) is
a connected one-dimensional subgroup of \(A_K\), the complement of \(E\) in (I8). It is not \(E\),
because \(i_E^\vee\lambda i_E\) is a polarization of \(E\), hence an isogeny. So it is \(F\), the only
other elliptic curve. Thus \(\pi_E=(\lambda^0_E)^{-1}i_E^\vee\lambda:A_K\to E\) is surjective with
kernel \(F\), and induces \(A_K/F\cong E\); in the same way \(\pi_F=(\lambda^0_F)^{-1}i_F^\vee\lambda\)
has kernel \(E\) and induces \(A_K/E\cong F\).

\emph{The gluing.} The polarization \(i_E^\vee\lambda i_E\) of \(E\) is \(n\lambda^0_E\) for some
\(n\ge1\), with kernel \(E[n]\); this kernel is also \(E\cap\ker(i_E^\vee\lambda)=E\cap F\). Likewise
\(E\cap F=F[n']\) with \(i_F^\vee\lambda i_F=n'\lambda^0_F\), and \(n=n'\) by counting. So
\(E\cap F=E[n]=F[n]\) inside \(A_K\). Let \(\varphi_n:E\times F\to A_K\) be the addition map. Its kernel
is \(\{(x,-x):x\in E\cap F\}\), the graph of \(\psi_n:E[n]\to F[n]\), \(x\mapsto-x\) (using
\(E[n]=F[n]\)). The pull-back \(\varphi_n^\vee\lambda\varphi_n\) is \(n\lambda_0\) with
\(\lambda_0=\lambda^0_E\times\lambda^0_F\), because \(i_E^\vee\lambda i_E=n\lambda^0_E\),
\(i_F^\vee\lambda i_F=n\lambda^0_F\) and \(i_F^\vee\lambda i_E=0\). Hence \(\ker\varphi_n\) is isotropic for
the pairing \((x,y)\mapsto e_n(x,\lambda_0y)\) on \((E\times F)[n]\), which is the product of the Weil
pairings of \(E\) and \(F\): if \(x,y\in\ker\varphi_n\) and \(y=ny'\), then \(\lambda_0y=\varphi_n^\vee\lambda\varphi_n(y')\)
with \(\varphi_n(y')\in A_K[n]\), and \(e_n(x,\varphi_n^\vee\lambda\varphi_ny')=e_n(\varphi_nx,\lambda\varphi_ny')=1\).
For \(x,y\in E[n]\) this reads \(e_n(x,y)\,e_n(\psi_nx,\psi_ny)=1\): \(\psi_n\) is an anti-isometry. Put
\(\Phi_n=(\pi_E,\pi_F):A_K\to E\times F\). For \(x\in E\) and \(y\in F\),
\(\Phi_n\varphi_n(x,y)=(\pi_Ex,\pi_Fy)=(nx,ny)\), so \(\Phi_n\varphi_n=[n]\), and
\(\varphi_n\Phi_n=[n]\) because \(\varphi_n\) is surjective. Since \(\lambda\) is defined over \(\Q\)
and \(\tilde\sigma(E)=F\), both maps commute with the action of \(G_{\Q}\) of (I7) on
\(E\times F=E\times E^\sigma\), so they descend to isogenies
\(\varphi_n:\Res_{K/\Q}E\to A\) and \(\Phi_n:A\to\Res_{K/\Q}E\) over \(\Q\).

\emph{The torsion of \(E(K)\).} Let \(G=\langle R_{45}\rangle\). Since \(\tilde\sigma\) fixes \(G\)
pointwise and maps \(E\) onto \(F\), the subgroups \(G\cap E\) and \(G\cap F\) of the cyclic group
\(G\) have the same order, hence coincide; call this subgroup \(H\) and put \(d=\#H\). Then \(E(K)\)
contains \(H\), of order \(d\), and \(\pi_E(G)\cong G/(G\cap F)=G/H\), of order \(45/d\). For
\(d\in\{1,5,9,45\}\) this gives a point of order \(45\) in \(E(K)\), contrary to (I1). Hence
\(d\in\{3,15\}\), \(E(K)\) has a point of order \(15\), and \(E(K)_{\mathrm{tors}}\cong\Z/15\) by
(I1); the same holds for \(F=E^\sigma\). This is (a).

\emph{The sign.} Put \(R_9=5R_{45}\). As \(d\in\{3,15\}\), the point \(P_0=3R_9\) of order \(3\) lies in
\(H\subseteq E\cap F=E[n]\); in particular \(3\mid n\). Let \(P\in E(K)\) be \(P_0\) regarded as a point
of \(E\). Since \(P_0\) is rational, \(P_0=\tilde\sigma(P_0)\), and \(\tilde\sigma(P_0)\) is the point
\(\sigma(P)\) of \(F=E^\sigma\). So \(\varphi_n(P,0)=P_0=\varphi_n(0,\sigma P)\), that is,
\((P,-\sigma P)\in\ker\varphi_n\), and as \(E(K)[3]=\langle P\rangle\),
\[\psi_n(Q)=-\sigma(Q)\qquad\text{for every }Q\in E(K)[3].\tag{1}\]
With \(+\sigma(Q)\) instead, condition (4) below would become \(a_3\sigma(a_3)\in K^{\times3}\);
\hyperref[rem-census]{Remark~3} tests the sign on census curves.

\emph{The \(3\)-adic valuation of \(n\).} Over \(K\),
\(\ker\Phi_n=\varphi_n(E[n]\times F[n])\cong(E[n]\times F[n])/\ker\varphi_n\), and
\((U,V)\mapsto V-\psi_n(U)\) identifies it \(G_K\)-equivariantly with \(F[n]\). Its \(K\)-rational
points form \(F(K)[n]\), which has no point of order \(9\). So \(\Phi_n(R_9)\neq0\); and
\(\Phi_n(R_9)\in(\Res_{K/\Q}E)(\Q)=E(K)\), where the \(3\)-primary torsion is \(\Z/3\). Hence
\(\Phi_n(R_9)=(P_1,\sigma P_1)\) with \(P_1\in E(K)\) of order \(3\). If \(9\mid n\), then
\(\varphi_n(P_1,\sigma P_1)=\varphi_n\Phi_n(R_9)=nR_9=0\), so \(\psi_n(P_1)=\sigma(P_1)\), against
(1). Therefore
\[v_3(n)=1.\tag{2}\]

\emph{The conjugate \((3,3)\)-gluing.} Write \(n=3m\) with \(3\nmid m\). The \(3\)-primary part
\(\Gamma_3\) of \(\ker\varphi_n\) is a characteristic subgroup of a subgroup defined over \(\Q\), hence
is defined over \(\Q\); by (2), over \(K\) it is the graph of \(\psi=\psi_n|_{E[3]}:E[3]\to F[3]\). Let
\(B=\Res_{K/\Q}(E)/\Gamma_3\) with quotient map \(\varphi\), and let \(\Phi:B\to\Res_{K/\Q}E\) be the
isogeny with \(\Phi\circ\varphi=[3]\), which exists because \(\Gamma_3\) is killed by \(3\). Both are
defined over \(\Q\). The map \(\psi\) is an anti-isometry: if \(m'm\equiv1\pmod3\), then
\(e_3(x,y)=e_n(m'x,y)\) for \(x,y\in E[3]\), likewise on \(F\), and \(\psi\) commutes with \(m'\). By (1),
\(\psi(P)=-\sigma(P)\), which is (b). As \(\Gamma_3\subseteq\ker\varphi_n\), we have
\(\varphi_n=\beta\circ\varphi\) for an isogeny \(\beta:B\to A\) over \(\Q\) whose kernel
\(\ker\varphi_n/\Gamma_3\) has exponent \(m\); so there is \(\beta':A\to B\) over \(\Q\) with
\(\beta\beta'=[m]\). The point \(R=\beta'(R_9)\in B(\Q)\) has order \(9\), since \(\beta(R)=mR_9\) has order
\(9\). Finally, over \(K\), \(\ker\Phi=\varphi(E[3]\times F[3])\cong(E[3]\times F[3])/\Gamma_3\) is
identified with \(F[3]\) as above, so its group of \(K\)-rational points is \(F(K)[3]\), of order \(3\).
As for \(R_9\), it follows that \(\Phi(R)=(P_1',\sigma P_1')\) with \(P_1'\in E(K)\) of order \(3\).
Replacing \(P\) by \(P_1'\), which also satisfies (b) by (1), gives (c). \qed

\medskip
Nothing later uses that \(B\) is principally polarized, or the description of \(B(\Q)_{\mathrm{tors}}\)
(it is \(\Z/45\): its \(3\)-primary part has order at most \(\#F(K)[3]\cdot\#E(K)[3]=9\), and its
prime-to-\(3\) part is that of \(E(K)\)). Section~\ref{conditions} uses only (a)--(c).

\section{Two necessary conditions}\label{conditions}

\begin{statement}{Lemma 4.}\label{lemma-conditions}
In the situation of \hyperref[lemma-parents]{Lemma~3}, choose a model
\(E:\ y^2+a_1xy+a_3y=x^3\) over \(K\) with \(P=(0,0)\). Then \(a_3\neq0\),
\[N_{K/\Q}(\Delta(E))\in\Q^{\times3},\tag{3}\]
\[a_3/\sigma(a_3)\in K^{\times3}.\tag{4}\]
Condition (3) uses only \hyperref[lemma-parents]{Lemma~3}(b).
\end{statement}

Such a model exists: \(P\) is a flex because \(3P=0\), and moving \(P\) to \((0,0)\) and its tangent
to \(y=0\) forces \(a_2=a_4=a_6=0\). The only changes of variables preserving this shape are
\((x,y)\mapsto(u^2x,u^3y)\), which multiply \(a_3\) by \(u^3\) and \(\Delta\) by \(u^{12}\); so (3) and
(4) do not depend on the model. The model of \(E^\sigma\) is the conjugate one, with \(\sigma(P)=(0,0)\).
By part~B of the certificate, \(\Delta(E)=a_3^3(a_1^3-27a_3)\), so \(a_3\neq0\); the
\(3\)-division polynomial is \(x\,h(x)\) with \(h(x)=3x^3+a_1^2x^2+3a_1a_3x+3a_3^2\); and
\(-P=(0,-a_3)\).

\emph{Proof of (3).} The extension
\(0\to\langle P\rangle\to E[3]\to E[3]/\langle P\rangle\to0\), with \(\langle P\rangle\cong\Z/3\) by
\(P\mapsto1\) and \(E[3]/\langle P\rangle\cong\mu_3\) by \(Q\mapsto e_3(P,Q)\), has a class
\(c(E)\) in \(\mathrm{Ext}^1_{G_K}(\mu_3,\Z/3)=H^1(K,\Hom(\mu_3,\Z/3))=H^1(K,\mu_3)=K^\times/K^{\times3}\);
here \(\Hom(\mu_3,\Z/3)\cong\mu_3\) canonically, as \(\F_3^\times=\{\pm1\}\). We use additive notation in
\(K^\times/K^{\times3}\).
\begin{enumerate}
\item[(a)] \(c(E)=\pm[\Delta(E)]\). Over \(K(\zeta_3)\), \(c(E)\) is the homomorphism describing the
unipotent action on \(E[3]\), so \(K(E[3])=K(\zeta_3,\sqrt[3]{c(E)})\). The roots of \(h\), the
\(x\)-coordinates of the points of \(E[3]\setminus\langle P\rangle\), are the numbers
\(x_k=-(e_k^2+a_1e_k+a_1^2)/9\), where \(e_k=\zeta_3^ke_0\) runs through the cube roots of
\(a_1^3-27a_3=\Delta/a_3^3\) (part~B of the certificate); they are distinct, since
\(x_0-x_1=-e_0(1-\zeta_3)\bigl(e_0(1+\zeta_3)+a_1\bigr)/9\) vanishes only if \(e_0^3\in\{0,a_1^3\}\),
that is, if \(\Delta=0\). Conversely
\(x_0+\zeta_3x_1+\zeta_3^2x_2=-e_0^2/3\), so \(e_0=e_0^3/e_0^2\) lies in \(K(\zeta_3,x_0,x_1,x_2)\).
So \(K(E[3])\) contains \(K(\zeta_3,\sqrt[3]\Delta)=K(\zeta_3,e_0)\), and it is generated over it by
\(y\)-coordinates, which satisfy quadratic equations; as \([K(E[3]):K(\zeta_3)]\) is \(1\) or \(3\), the two fields coincide. By
Kummer theory \(c(E)\) and \([\Delta]\) generate the same subgroup of
\(K(\zeta_3)^\times/K(\zeta_3)^{\times3}\), hence of \(K^\times/K^{\times3}\), the restriction map
being injective as \([K(\zeta_3):K]\le2\).
\item[(b)] \(c(E^\sigma)=\sigma(c(E))\), by transport of structure under \(\tilde\sigma\), which maps
\(E\), \(P\), \(E[3]\) and \(e_3\) to \(E^\sigma\), \(\sigma(P)\), \(E^\sigma[3]\) and \(e_3\); and
\(\Delta(E^\sigma)=\sigma(\Delta(E))\).
\item[(c)] \(c(E^\sigma)=-c(E)\). By (1), \(\psi\) maps \(\langle P\rangle\) to \(\langle\sigma P\rangle\)
by \(\alpha=-1\); it induces on the quotients, identified with \(\mu_3\), multiplication by the exponent
\(\beta\) with \(e_3(\sigma P,\psi Q)=e_3(P,Q)^\beta\). As \(\psi\) is an anti-isometry,
\(e_3(P,Q)^{-1}=e_3(\psi P,\psi Q)=e_3(-\sigma P,\psi Q)=e_3(P,Q)^{-\beta}\), so
\[\beta=+1.\]
An isomorphism of extensions that is \(\alpha\) on the subgroups and \(\beta\) on the quotients gives
\(\alpha\,c(E)=\beta\,c(E^\sigma)\), so \(c(E^\sigma)=-c(E)\).
\item[(d)] By (a)--(c), \(c(E)+\sigma(c(E))=0\), that is, \(N_{K/\Q}(\Delta)^{\pm1}\in K^{\times3}\).
Finally \(K^{\times3}\cap\Q^\times=\Q^{\times3}\): if \(q=k^3\) with \(k\in K\), then
\(q^2=N_{K/\Q}(k)^3\), so \(q=q^3/q^2\) is a rational cube. This proves (3).
\end{enumerate}

\emph{Proof of (4).} Let \(R\), \(\varphi\), \(\Phi\) be as in \hyperref[lemma-parents]{Lemma~3}(c),
with \(\Phi(R)=(P,\sigma P)\). For \(U\in E(\bar\Q)\) with \(3U=P\) and \(V\in E^\sigma(\bar\Q)\) with
\(3V=\sigma P\) we have \(\Phi\varphi(U,V)=(3U,3V)\). So the fibre
\(\Phi^{-1}(P,\sigma P)=\varphi([3]^{-1}P\times[3]^{-1}\sigma P)\) is a torsor under
\(\ker\Phi=\varphi(E[3]\times E^\sigma[3])\) with the \(K\)-rational point \(R\). Transport it by the
\(G_K\)-isomorphism \(\theta:\ker\Phi\to E^\sigma[3]\), \(\varphi(U',V')\mapsto V'-\psi(U')\), which is
well defined because \(\ker\varphi\) is the graph of \(\psi\). If \(\kappa\) denotes the Kummer map
\(E(K)\to H^1(K,E[3])\) of multiplication by \(3\), the class of the transported torsor is
\(\kappa(\sigma P)-\psi_*\kappa(P)\in H^1(K,E^\sigma[3])\). As the torsor has a \(K\)-rational point,
\[\kappa(\sigma P)=\psi_*\kappa(P).\]
Let \(\lambda_E:E(K)\to H^1(K,\mu_3)=K^\times/K^{\times3}\) be \(\kappa\) followed by
\(E[3]\to E[3]/\langle P\rangle\cong\mu_3\), and define \(\lambda_{E^\sigma}\) in the same way with
\(\sigma P\). Since \(\psi\) induces \(\beta=+1\) on the quotients,
\(\lambda_{E^\sigma}(\sigma P)=\lambda_E(P)\); by transport of structure,
\(\lambda_{E^\sigma}(\sigma P)=\sigma(\lambda_E(P))\). So \(\lambda_E(P)\) is fixed by \(\sigma\).

To compute \(\lambda_E(P)\), let \(\phi:E\to E'=E/\langle P\rangle\) be the quotient isogeny and
\(\hat\phi:E'\to E\) its dual, so that \(\hat\phi\phi=[3]\). As \(\phi\) maps \([3]^{-1}(Q)\) onto
\(\hat\phi^{-1}(Q)\) and \(E[3]/\langle P\rangle\) isomorphically onto \(\ker\hat\phi\), \(\lambda_E\) is
the connecting map \(\delta\) of \(0\to\ker\hat\phi\to E'\to E\to0\), followed by an isomorphism
\(\ker\hat\phi\cong\mu_3\). Part~B of the certificate checks that
\(E':\ y^2+a_1xy+a_3y=x^3-5a_1a_3x-a_1^3a_3-7a_3^2\), that \(\hat\phi\phi=[3]\), and that
\[\hat\phi(x,y)=\Bigl(\frac{x^3-9a_1a_3x-a_1^3a_3-27a_3^2}{(3x+a_1^2)^2},\
\Bigl(\frac{y-4a_3}{3x+a_1^2}\Bigr)^3\Bigr).\]
Thus \(y\circ\hat\phi=w^3\) with \(w=(y-4a_3)/(3x+a_1^2)\) in the function field of \(E'\). For
\(S\in\ker\hat\phi\) the function \(w(\cdot+S)/w\) has cube \(1\), hence is a constant \(\xi(S)\in\mu_3\),
and \(\xi:\ker\hat\phi\to\mu_3\) is a homomorphism; it is injective, since otherwise
\(w=w_0\circ\hat\phi\) with \(w_0^3=y\), which is impossible because \(\operatorname{div}(y)=3(P)-3(O)\)
and \((P)-(O)\) is not principal. For \(Q\in E(K)\setminus\{O,P\}\) and \(Q'\) with \(\hat\phi(Q')=Q\) we
get \(w(Q')^3=y(Q)\in K^\times\) and \(g(w(Q'))/w(Q')=\xi(gQ'-Q')\) for \(g\in G_K\), so
\(\xi_*\delta(Q)=[y(Q)]\). With \(Q=-P\) and \(y(-P)=-a_3\) this gives \(\lambda_E(-P)=\pm[a_3]\) (two
isomorphisms \(\ker\hat\phi\cong\mu_3\) differ by a sign, and \(-1\) is a cube), and
\(\lambda_E(P)=-\lambda_E(-P)=\mp[a_3]\). As \(\lambda_E(P)\) is fixed by \(\sigma\), so is \([a_3]\),
which is (4). \qed

\medskip
The sign \(\beta=+1\) is where the minus sign in (1) enters: with \(\beta=-1\) the argument would give
\(N_{K/\Q}(a_3)\in K^{\times3}\) instead of (4).

Put \(S=\operatorname{Tr}_{K/\Q}(a_3)\) and \(N=N_{K/\Q}(a_3)\neq0\). If \(w\in K\) satisfies
\(w^3=a_3/\sigma(a_3)\), then \(N_{K/\Q}(w)^3=1\), so \(N_{K/\Q}(w)=1\) as it is rational, and
\(r=w+\sigma(w)\in\Q\) satisfies, since \(w\sigma(w)=1\),
\[r^3-3r=w^3+\sigma(w)^3=\frac{a_3}{\sigma(a_3)}+\frac{\sigma(a_3)}{a_3}=\frac{S^2}{N}-2.\tag{5}\]
If \(v_\ell(r^3-3r)<0\) for a prime \(\ell\), then \(v_\ell(r)<0\), so \(v_\ell(r^3)=3v_\ell(r)<v_\ell(3r)\)
and
\[v_\ell(r^3-3r)=3v_\ell(r)\equiv0\pmod3.\tag{6}\]

\section{A primitive cyclotomic equation}\label{cyclotomic}

\begin{statement}{Proposition 5.}\label{prop-cyclo}
Let \(a\), \(b\) be coprime integers with
\(\Phi_{10}(a,b)=a^4-a^3b+a^2b^2-ab^3+b^4\in\{c^3,5c^3\}\) for some integer \(c\). Then
\((a:b)\) is one of \((0:1)\), \((1:0)\), \((1:1)\), \((-1:1)\).
\end{statement}

\emph{Proof.} The descent follows Poonen \cite[Section~4]{Po}, who applies it to \(x+\zeta y\) for a
solution of \(x^5+y^5=z^3\) (\cite[(15), p.~198, and p.~199]{Po}). Let \(L=\Q(\zeta)\) with
\(\zeta=\zeta_5\), and \(\tau=-(\zeta^2+\zeta^3)=(1+\sqrt5)/2\). The ring of integers of \(L\) is
\(\Z[\zeta]\), and its Minkowski bound \((4!/4^4)(4/\pi)^2\sqrt{125}\approx1.699\) is less than \(2\),
so \(L\) has class number one. Its unit group is \(\mu_{10}\times\tau^{\Z}\): for a unit \(u\),
\(u/\bar u\) is an algebraic integer all of whose conjugates have absolute value \(1\), hence a root
of unity, and it is \(\equiv1\) modulo \(\pi=1-\zeta\), so it lies in \(\mu_5\); choosing
\(\eta\in\mu_5\) with \(\eta^2=u/\bar u\), the unit \(u/\eta\) is real, hence \(\pm\tau^k\).

Put \(\alpha=a+b\zeta\). Then \(N_{L/\Q}(\alpha)=\Phi_{10}(a,b)\), a positive integer. If a prime
\(\mathfrak p\) divides \(\alpha\) and a conjugate \(a+b\zeta^k\) with \(k\neq1\), then it divides
\(b(\zeta-\zeta^k)\), which is \(b\pi\) times a unit; if \(\mathfrak p\neq(\pi)\), it divides \(b\), hence
\(a\), which is impossible. So the conjugates of \(\alpha\) are pairwise coprime away from \(\pi\).
Now let \(\mathfrak p\nmid\pi\) divide \(\alpha\), above the prime \(\ell\neq5\). No other conjugate of
\(\mathfrak p\) divides \(\alpha\): if \(g\mathfrak p\mid\alpha\) with \(g\mathfrak p\neq\mathfrak p\), then
\(\mathfrak p\) divides both \(\alpha\) and \(g^{-1}\alpha\). For the same reason no
\(g\neq1\) fixes \(\mathfrak p\), so \(\ell\) splits completely, the four primes above \(\ell\) are the
conjugates of \(\mathfrak p\), and \(v_{\mathfrak p}(\alpha)=v_\ell(\Phi_{10}(a,b))\equiv0\pmod3\). At
\(\pi\): \(\alpha=(a+b)-b\pi\), so \(v_\pi(\alpha)=0\) if \(5\nmid a+b\), and \(v_\pi(\alpha)=1\) if
\(5\mid a+b\) (then \(5\nmid b\) and \(v_\pi(a+b)\ge v_\pi(5)=4\)). By class number one,
\(\alpha=\varepsilon\pi^iw_0^3\) with \(\varepsilon\) a unit, \(i\in\{0,1\}\) and \(w_0\in L\). Roots of
unity are cubes, since \(\gcd(3,10)=1\), and
\[\pi^4=5\zeta^2\tau^{-2},\qquad\pi=5\tau\Bigl(\frac{\zeta^4}{\tau\pi}\Bigr)^3,\]
so
\[\alpha=5^i\tau^jw^3\qquad\text{with }i\in\{0,1\},\ j\in\{0,1,2\},\ w\in L^\times.\tag{7}\]

\emph{Modulo \(2\).} The prime \(2\) is inert in \(L\), with residue field \(\F_{16}\), whose
multiplicative group is cyclic of order \(15\); its cubes are its elements of order dividing \(5\). As
\(a\), \(b\) are coprime, \(\alpha\) is a \(2\)-adic unit, hence so is \(w\). The reduction
\(\bar\tau\) is a root of \(X^2+X+1\), of order \(3\), so it is not a cube, and \(\bar5=1\). If exactly
one of \(a\), \(b\) is odd, then \(\bar\alpha\in\{1,\bar\zeta\}\) is a cube (\(\bar\zeta^5=1\)), so
\(j=0\). If both are odd, then \(\bar\alpha=1+\bar\zeta\), and
\((1+\bar\zeta)^5=(1+\bar\zeta^4)(1+\bar\zeta)=\bar\zeta+\bar\zeta^4=\bar\tau^2=\bar\tau^5\), so
\(1+\bar\zeta\) is \(\bar\tau\) times a cube and \(j=1\). Thus \(j\in\{0,1\}\).

\emph{Two curves over \(L\).} The automorphisms \(\zeta\mapsto\zeta^2\) and \(\zeta\mapsto\zeta^3\) map
\(\tau\) to \(-\tau^{-1}\). Multiplying (7) and its images under these two automorphisms gives
\((a+b\zeta)(a+b\zeta^2)(a+b\zeta^3)=5^{3i}\tau^{-j}W^3\) with \(W\in L\). If \(b\neq0\), dividing by
\(b^3\) shows that \(t=a/b\) is the first coordinate of an \(L\)-point of
\[D_j:\quad Y^3=\tau^j(t+\zeta)(t+\zeta^2)(t+\zeta^3),\qquad j\in\{0,1\};\tag{8}\]
the factor \(5^{3i}\) is a cube. The case \(b=0\), that is \((a:b)=(1:0)\) with \(j=0\), is the point \((t:Y:W)=(1:1:0)\) at
infinity of the projective model \(Y^3=\tau^j(t+\zeta W)(t+\zeta^2W)(t+\zeta^3W)\). Each \(D_j\) is a smooth plane cubic with the
\(L\)-rational flex \((-\zeta,0)\), hence an elliptic curve. With this point as origin, \(D_j\) is
isomorphic over \(L\) to
\[\begin{aligned}
E_0&:\ y^2+(4\zeta^3+3\zeta^2+7\zeta+6)y=x^3-10\zeta^3-5\zeta^2-10\zeta,\\
E_1&:\ y^2+(\zeta^3+7\zeta^2+13\zeta+9)y=x^3-35\zeta^3-55\zeta^2-35\zeta.
\end{aligned}\]
Part~C of the certificate proves: \(N_{L/\Q}(\Delta(E_j))=3^{12}5^{12}\); \(E_j\) has rank \(1\) and
\(E_j(L)\cong\Z/3\oplus\Z\) with proved generators (Magma's \texttt{RankBounds} and
\texttt{MordellWeilGroup}, with class groups computed unconditionally up to the Minkowski bound); and elliptic Chabauty \cite{Br} at the prime \(11\), which splits
completely in \(L\) and is a prime of good reduction, returns an upper bound equal to the number of
points found and no unresolved residue classes. The saturation it requires (at \(2\) for \(j=0\), none
for \(j=1\)) holds because the Mordell--Weil groups are proved to be complete. The \(L\)-points of
\(D_j\) with rational \(t\)-coordinate are therefore exactly those with
\[t\in\{0,\infty\}\ \text{ on } D_0,\qquad t\in\{1,-1\}\ \text{ on } D_1,\tag{9}\]
which gives the four values of \((a:b)\). \qed

\medskip
In \cite[Section~4]{Po} the descent (7) is followed by an analysis of the curves
\((U+1)(U^5+1)=\eta V^3\) and their elliptic quotients over \(\Q(\sqrt5)\). Here the equation is
different (the norm \(\Phi_{10}(a,b)\) itself is a cube or five times a cube); what is new is that the
product of three conjugates removes the factor \(5^{3i}\), so that only the two cubic twists
\(D_0\), \(D_1\) remain, and these are treated by elliptic Chabauty over \(\Q(\zeta_5)\).

\section{The quadratic points of \(X_1(15)\)}\label{families}

\begin{statement}{Proposition 6.}\label{prop-families}
Let \(K\) be a quadratic field, \(E\) an elliptic curve over \(K\) with
\(E(K)_{\mathrm{tors}}\cong\Z/15\), and \(P\in E(K)\) a point of order \(3\). Write
\(E:\ y^2+a_1xy+a_3y=x^3\) with \(P=(0,0)\). Then (3) and (4) do not both hold.
\end{statement}

\emph{A model of \(X_1(15)\).} Let \(P_5\in E(K)\) have order \(5\). By (I9),
\((E,P_5)\cong(E_d,(0,0))\) over \(K\) for some \(d\in K\) with \(d^5(d^2-11d-1)\neq0\); let
\((x_3,y_3)\) be the image of \(P\). The \(3\)-division polynomial of \(E_d\) is
\[\psi_3(X)=3X^4+(d^2-6d+1)X^3+3(d^2-d)X^2+3d^2X-d^3,\]
and \(\psi_3(0)=-d^3\neq0\), \(\psi_3(d)=d^5\neq0\). So \(u=x_3/d\) is neither \(0\) nor \(1\), and with
\(z=du/(u-1)\) one has \(\psi_3(du)/d^3=(u-1)^2\bigl(uz^2+3u(u-1)z+(u-1)\bigr)\) (part~D of the
certificate). Hence \((u,z,y_3)\) is a \(K\)-point of the affine curve \(X^{\mathrm{aff}}\) given by
\[\begin{aligned}
&uz^2+3u(u-1)z+(u-1)=0,\qquad d=\frac{(u-1)z}{u},\qquad x_3=(u-1)z,\\
&y_3^2+(1-d)x_3y_3-dy_3=x_3^3-dx_3^2,
\end{aligned}\tag{10}\]
where the last equation, multiplied by \(u^2\), is a polynomial equation in \(u,z,y_3\). The point
\((0,0)\) of \(E_d\) has order \(5\) and \((x_3,y_3)\) has order \(3\). Translating \((x_3,y_3)\) to
\((0,0)\) and shearing its tangent to \(y=0\) turns \(E_d\) into a model as in
\hyperref[prop-families]{Proposition~6}, with
\[a_3=2y_3+(1-d)x_3-d,\qquad \delta=d^2(d^2-11d-1);\tag{11}\]
the translation contributes \(2y_3+(1-d)x_3\) to the coefficient of \(y\) and the shear does not change
it, and the discriminant \(d^5(d^2-11d-1)\) is unchanged and differs from \(\delta\) by the cube
\(d^3\). So, by the remarks after \hyperref[lemma-conditions]{Lemma~4}, (3) and (4) are conditions on
\(\delta\) and \(a_3\) at the point \((u,z,y_3)\). We call a point of (the projective closure of)
\(X^{\mathrm{aff}}\) a \emph{cusp} if \(d\) is \(0\) or \(\infty\) or a root of \(d^2-11d-1\) there;
the point \((u,z,y_3)\) is not a cusp. The curve \(X^{\mathrm{aff}}\) is irreducible, and its singular
points are \((1,0,0)\), where \(d=0\), and two conjugate points with \(d^2-11d-1=0\) (part~D). So \((u,z,y_3)\) is a smooth point of \(X^{\mathrm{aff}}\).

\emph{The curve \(W\).} Let \(W:\ y^2+xy+y=x^3+x^2\) (Cremona \(15a8\)), with origin \(O\) and the
points \(T_1=(0,0)\), \(T_2=(-1,0)=2T_1\), \(T_3=(0,-1)=3T_1\). It has rank \(0\) and
\(W(\Q)=\{O,T_1,T_2,T_3\}\cong\Z/4\) (part~D). The rational functions
\[\begin{aligned}
u&=\frac{x+1}{x^2+x+1},\qquad z=\frac{x\,(x^2-1-(x+2)y)}{(x+1)(x^2+x+1)},\\
y_3&=\frac{x^5\bigl(x-2x^2-2x^3+(x^2+4x+1)y\bigr)}{(x+1)^2(x^2+x+1)^3}
\end{aligned}\]
on \(W\) define a birational map \(W\to X^{\mathrm{aff}}\), with inverse
\[\begin{aligned}
x&=\frac{(u-1)^3(u+1)^2+u^2(u-1)^2z+u^2(-9u^3+3u^2+7u+1)y_3-u^3(3u+2)zy_3}{u(u-1)^2(u^2-u-1)},\\
y&=-\frac{(u-1)^4+(u-1)^3z+3u^2(u-1)y_3+u^2zy_3}{u(u-1)^3}
\end{aligned}\]
(part~D checks that the image lies on (10) and that the composite \(W\to X^{\mathrm{aff}}\to W\) is the
identity). On \(W\),
\[d=\frac{x^3\bigl(1-x^2+(x+2)y\bigr)}{(x+1)^2(x^2+x+1)},\qquad x_3=du;\]
\(d\) has poles at \(O\) and \(T_2\) and zeros at \(T_1\) and \(T_3\), so the four rational points of
\(W\) are cusps. As \((u,z,y_3)\) is a smooth point of \(X^{\mathrm{aff}}\), there is a unique point
\(Q\in W(K)\) at which the functions \(u\), \(z\), \(y_3\) take these values, and the functions
\(a_3\) and \(\delta\) of (11) are regular at \(Q\). Since \(Q\) is not a cusp, \(Q\notin W(\Q)\), so
\(Q\) is a quadratic point with residue field \(K\), with conjugate \(Q^\sigma\).

\emph{The four families.} The divisor \((Q)+(Q^\sigma)\) is rational of degree \(2\), so it is
linearly equivalent to \((O)+(T)\) with \(T=Q+Q^\sigma\in W(\Q)\). For each \(T\in W(\Q)\), by
Riemann--Roch, the effective divisors linearly equivalent to \((O)+(T)\) are the fibres of a morphism
\(W\to\mathbb P^1\) of degree \(2\), unique up to an automorphism of \(\mathbb P^1\); we fix it as the
function \(t_j\) below, whose divisor of poles lies in this class. So \((Q)+(Q^\sigma)\) is the fibre
of some \(t_j\) over a point \(t_0\in\mathbb P^1(\Q)\); as \(Q\) is not a cusp and the poles of
\(t_j\) are rational points, \(t_0\in\Q\).
\begin{enumerate}
\item[(1)] The class of \((O)+(T_2)\): \(t_1=y/(x+y+1)\), with divisor \(2(T_1)-2(T_3)\) and
\(t_1(O)=t_1(T_2)=1\). Over \(\Q(t_1)\), \((1-t_1)^2x^2-t_1x-t_1=0\) and \(y=t_1(x+1)/(1-t_1)\).
\item[(2)] The class of \(2(O)\): \(t_2=x+1\), with divisor \(2(T_2)-2(O)\) and
\(t_2(T_1)=t_2(T_3)=1\). Over \(\Q(t_2)\), \(y^2+t_2y-t_2(t_2-1)^2=0\) and \(x=t_2-1\).
\item[(3)] The class of \((O)+(T_3)\): \(t_3=y/x\), with divisor \((T_1)+(T_2)-(T_3)-(O)\),
normalized by \(t_3=-1\) at the cusp \((\omega,-\omega)\) of degree \(2\), where \(\omega^2+\omega+1=0\).
Over \(\Q(t_3)\), \(x^2+(1-t_3-t_3^2)x-t_3=0\) and \(y=t_3x\).
\item[(4)] The class of \((O)+(T_1)\): \(t_4=-(x+y+1)/x\), with divisor \((T_2)+(T_3)-(T_1)-(O)\),
normalized by \(t_4=-1\) at the cusp \((\omega,-1)\). Over \(\Q(t_4)\), \(x^2+(1-t_4-t_4^2)x-t_4=0\)
and \(y=-(t_4+1)x-1\).
\end{enumerate}
Each \(t_j\) is determined by its divisor and its normalization, so the families do not depend on
any choice (part~D checks the divisors and values). The negation of \(W\) exchanges families \(3\)
and \(4\) and maps \(t_1\) to \(1/t_1\).

\emph{Traces and norms.} In each family the function field \(\Q(W)\) is the quadratic extension of
\(\Q(t)\), \(t=t_j\), given above. Computing in it (part~D), the function
\(R_j=\operatorname{Tr}(a_3)^2/N(a_3)-2\in\Q(t)\), formed with the trace and norm from \(\Q(W)\) to
\(\Q(t)\), is \(R_j=N_j(t)/D_j(t)\), where, with \(B(t)=t^2-t+1\), \(\Phi_{10}(t)=t^4-t^3+t^2-t+1\)
and \(\Phi_5(t)=t^4+t^3+t^2+t+1\),
\[D_1=D_2=B^3\Phi_{10},\qquad D_3=D_4=t^9\Phi_5,\]
\[\begin{aligned}
N_1&=-2t^{10}-17t^9+60t^8-30t^7-126t^6+228t^5-126t^4-30t^3+60t^2-17t-2,\\
N_2&=2t^{10}+8t^9-24t^8+30t^7-45t^6+60t^5-45t^4+30t^3-24t^2+8t+2,\\
N_3&=-t^{22}-6t^{21}-7t^{20}+16t^{19}+21t^{18}-44t^{17}-37t^{16}+88t^{15}+38t^{14}-132t^{13}-17t^{12}\\
&\quad+152t^{11}-17t^{10}-132t^9+38t^8+88t^7-37t^6-44t^5+21t^4+16t^3-7t^2-6t-1,\\
N_4&=-t^{22}-5t^{21}-2t^{20}+15t^{19}-8t^{18}-37t^{17}+55t^{16}+20t^{15}-119t^{14}+89t^{13}+73t^{12}\\
&\quad-170t^{11}+73t^{10}+89t^9-119t^8+20t^7+55t^6-37t^5-8t^4+15t^3-2t^2-5t-1.
\end{aligned}\]
The \(N_j\) have integer coefficients, are palindromic, and
\[\Res(\Phi_{10},N_1)=\Res(\Phi_{10},N_2)=\Res(\Phi_5,N_3)=\Res(\Phi_5,N_4)=25.\]
Moreover the norm of \(\delta\) from \(\Q(W)\) to \(\Q(t)\) is \(F_j(t)\,g_j(t)^3\) with
\(g_j\in\Q(t)\) (computed in part~D) and
\[\begin{aligned}
F_1=F_2&=(t^2-t+1)(t^4-11t^3+21t^2-11t+1),\\
F_3=F_4&=(t^2-t+1)^2(t+1)(t^4+t^3-9t^2+t+1).
\end{aligned}\]

\emph{Specialization.} Let \(\mathcal O\) be the integral closure of the local ring
\(\Q[t]_{(t-t_0)}\) in \(\Q(W)\). As \(\{Q,Q^\sigma\}\) is the whole fibre of \(t\) over \(t_0\),
\(\mathcal O\) is the local ring of \(W\) at this point of degree \(2\); it is free of rank \(2\) over
\(\Q[t]_{(t-t_0)}\), with \(\mathcal O/(t-t_0)\mathcal O\cong K\), and \(a_3,\delta\in\mathcal O\). Traces
and norms are traces and determinants of multiplication maps in a basis, so they commute with
reduction modulo \(t-t_0\): with \(S\) and \(N\) as in (5), \(S^2/N-2=R_j(t_0)\), and
\(N_{K/\Q}(\delta(Q))=F_j(t_0)g_j(t_0)^3\) whenever \(g_j(t_0)\) is finite and nonzero.

\emph{The arithmetic.} Suppose that (4) holds. By (5), \(r^3-3r=R_j(t_0)\) for some \(r\in\Q\). Write
\(t_0=a/b\) with coprime integers \(a\), \(b\), \(b>0\). In families \(1\) and \(2\),
\(R_j(t_0)=N_j(a,b)/\bigl(B(a,b)^3\Phi_{10}(a,b)\bigr)\), a quotient of binary forms of degree \(10\).
Let \(\ell\neq5\) be a prime dividing \(\Phi_{10}(a,b)\). Then \(\ell\nmid b\) (otherwise \(\ell\mid a\)),
and \(\ell\nmid N_j(a,b)\): the resultant of the binary forms \(\Phi_{10}\) and \(N_j\) is
\(\Res(\Phi_{10},N_j)=25\), as both have full degree in \(t\), and it times a power of \(a\) and times
a power of \(b\) lie in the ideal generated by \(\Phi_{10}(a,b)\) and \(N_j(a,b)\). So \(v_\ell(R_j(t_0))=-3v_\ell(B(a,b))-v_\ell(\Phi_{10}(a,b))<0\), and (6) gives
\(v_\ell(\Phi_{10}(a,b))\equiv0\pmod3\). Also \(\Phi_{10}(a,b)=N_{L/\Q}(a+b\zeta)>0\) and
\(v_5(\Phi_{10}(a,b))=v_\pi(a+b\zeta)\in\{0,1\}\) (Section~\ref{cyclotomic}). Hence \(\Phi_{10}(a,b)\) is
a cube or five times a cube, and \hyperref[prop-cyclo]{Proposition~5} gives
\(t_0\in\{0,\infty,1,-1\}\). In families \(3\) and \(4\), if \(t_0\neq0\), then
\(R_j(t_0)=N_j(a,b)/\bigl(a^9b^9\Phi_5(a,b)\bigr)\), a quotient of forms of degree \(22\); a prime
\(\ell\neq5\) dividing \(\Phi_5(a,b)\) divides neither \(a\), \(b\) nor \(N_j(a,b)\), so in the same way
\(\Phi_5(a,b)=\Phi_{10}(-a,b)\) is a cube or five times a cube, and Proposition~5 for \((-a,b)\) gives
the same four values of \(t_0\).

\emph{The four fibres.} Part~D determines the fibres of \(t_j\) over \(0\), \(1\), \(-1\) and
\(\infty\). Three of them consist of cusps, and the fourth, over \(t_0\), is one non-cuspidal point
with residue field \(K\), where \(d\) is rational:
\[\begin{array}{c|c|c|c|c|c}
j&\text{fibres of cusps}&t_0&K&d&N_{K/\Q}(\Delta)\bmod\Q^{\times3}\\\hline
1&0,\ 1,\ \infty&-1&\Q(\sqrt{-15})&-\tfrac18&F_1(-1)=135\\
2&0,\ 1,\ \infty&-1&\Q(\sqrt{-15})&8&F_2(-1)=135\\
3&0,\ -1,\ \infty&1&\Q(\sqrt5)&\tfrac12&F_3(1)=-10\\
4&0,\ -1,\ \infty&1&\Q(\sqrt5)&-2&F_4(1)=-10
\end{array}\]
At the non-cuspidal point over \(t_0\) the curve \(E_d\) has \(E_d(K)_{\mathrm{tors}}\cong\Z/15\), the
cube cofactor \(g_j(t_0)\) is finite and nonzero (it is \(5/768\), \(80/3\), \(-5/8\), \(-10\) in
families \(1\)--\(4\)), and \(N_{K/\Q}(\Delta(E_d))\) is \(135\) resp.~\(-10\) times a rational cube.
Neither \(135\) nor \(-10\) is a rational cube, so (3) fails there. Hence (3) and (4) never hold
together. \qed

\section{Proof of the Theorem and consequences for the census}\label{census}

\emph{Proof of the Theorem.} Let \(A/\Q\) be principally polarized, simple over \(\Q\) and
geometrically split, with a rational point of order \(45\). By (I3), \(A\) is of type (N), (CM) or
(S). Types (CM) and (S) are excluded by \hyperref[prop-A]{Proposition~A}. In type (N),
\hyperref[lemma-parents]{Lemma~3} gives a quadratic field \(K\), an elliptic curve \(E\) over \(K\)
with \(E(K)_{\mathrm{tors}}\cong\Z/15\) and a point \(P\in E(K)\) of order \(3\), and
\hyperref[lemma-conditions]{Lemma~4} shows that (3) and (4) hold for the model of \(E\) with
\(P=(0,0)\). This contradicts \hyperref[prop-families]{Proposition~6}. \qed

\medskip
\emph{Proof of the Corollary.} The Jacobian of a curve of genus \(2\) over \(\Q\) is a principally
polarized abelian surface, and a group containing an element of order \(45\) contains a cyclic
subgroup of order \(45\). \qed

\medskip
Among the groups listed on the census, \([45]\) is the only one with an element of order \(45\). It
becomes impossible in the class ``geometrically split, simple over \(\Q\)''. Simplicity over \(\Q\)
is essential: the class ``split over \(\Q\)'' contains the curve 45.a and a certified infinite family
of Jacobians with exact torsion group \([45]\) (see the
\href{https://genus-2-torsion.github.io/group.html?g=45}{census page}), so the combined
geometrically split cell of \([45]\) keeps its grade. The class ``geometrically simple'' is not
affected by this proof.

\section{Remarks}\label{remarks}

\begin{enumerate}
\item\label{rem-hyp} \emph{The hypotheses.} The principal polarization and the simplicity over \(\Q\)
are used only in Section~\ref{parents}. Simplicity makes the two elliptic curves conjugate. The
principal polarization identifies \(A_K/F\) with \(E\) itself, so that the torsion of \(G/H\) lands in
\(E(K)\), and makes \(A_K\) an \((n,n)\)-gluing with anti-isometric kernel; for a polarization of
higher degree \(A_K/F\) is only isogenous to \(E\). Proposition~A uses neither hypothesis;
Lemma~2 and Propositions~5 and~6 are statements about elliptic curves and a Diophantine equation;
and Lemma~4 uses only the output (a)--(c) of Lemma~3.

\item\label{rem-both} \emph{Both conditions are needed.} At the non-cuspidal points of the table in
Section~\ref{families}, condition (4) holds (\(a_3/\sigma(a_3)=-1\), so \(r=-2\) and
\(R_j(t_0)=-2\)), and only (3) excludes them. There \(d\) is rational: the curves are the base changes
to \(\Q(\sqrt{-15})\) of the curve \(50b2\) (\(d=-\tfrac18\), \(8\)) and to \(\Q(\sqrt5)\) of the curve
\(50b1\) (\(d=\tfrac12\), \(-2\)), so \(E\cong E^\sigma\). Conversely, at the point of family \(1\)
over \(t_1=2\), over \(\Q(\sqrt3)\), condition (3) holds and (4) fails; this point gives the elliptic
curve of the census curve 15.c (\hyperref[rem-census]{Remark~3}).

\item\label{rem-census} \emph{A consistency check with census curves (not used in the proof).} The
census curves \href{https://genus-2-torsion.github.io/curve.html?id=9.c}{9.c},
\href{https://genus-2-torsion.github.io/curve.html?id=18.c}{18.c} and
\href{https://genus-2-torsion.github.io/curve.html?id=15.c}{15.c} are geometrically split and simple
over \(\Q\), with rational torsion subgroups \([9]\), \([18]\) and \([15]\). The census records for each
a map of degree \(3\) to an elliptic curve \(E\) over \(K=\Q(i)\), \(\Q(\sqrt6)\) and \(\Q(\sqrt3)\).
As \(3\) is prime, the map does not factor through an isogeny, so the induced homomorphism
\(E\to J_K\) is injective \cite[Proposition~11.4.3]{BL}. Part~E of the certificate checks that \(E\)
has no complex multiplication and is not geometrically isogenous to \(E^\sigma\) (the squares of the
traces of Frobenius at two conjugate primes differ), so these Jacobians are of type (N) and \(E\) is one
of the two elliptic curves of Section~\ref{parents}; and that \(E(K)_{\mathrm{tors}}\) is \(\Z/3\),
\(\Z/6\) and \(\Z/15\). For 9.c and
18.c the argument of Section~\ref{parents}, with the subgroup of order \(9\) in place of the group of
order \(45\), forces \(d=3\) and hence (1), (3) and (4); and indeed (3) and (4) hold for both. For
15.c, which has no rational point of order \(9\), (3) holds and (4) fails. The variant
\(a_3\sigma(a_3)\in K^{\times3}\) of (4) with the opposite sign in (1) fails for all three curves.

\item\label{rem-rerun} \emph{What was not used.} The proof uses no bound on the height of points, no
search for points as a completeness argument and no conjecture: the ranks, the Mordell--Weil groups
and the Chabauty bounds of Section~\ref{cyclotomic} and the rank of \(W\) are proved.
\end{enumerate}

\section{Computations}\label{computations}

The certificate
\href{https://genus-2-torsion.github.io/data/knowledge/sources/qsimple_45_checks.m}{\texttt{qsimple\_45\_checks.m}}
(Magma V2.30-1, about \(11\) seconds) consists of assertions and ends with
\texttt{ALL\_CHECKS\_PASSED}; a failed assertion stops it before that line. Its log is
\href{https://genus-2-torsion.github.io/data/knowledge/sources/qsimple_45_checks.log}{\texttt{qsimple\_45\_checks.log}}.
SHA-256 of the script:
\texttt{3898e5d57d639cd4\allowbreak 3555feb97c996164\allowbreak ef867549fd76b5df\allowbreak 02bf40c2da0a0fe6};
of the log:
\texttt{4681daf8118e304a\allowbreak 60d91375822b1a88\allowbreak 74028e922ed94e4e\allowbreak a63803f907b3f060}.
Part~A checks the finite statements of (I1), the residues modulo \(5\) and \(3\) used in
Section~\ref{square} and \((1+\sqrt2)^4<34\). Part~B checks, for \(y^2+a_1xy+a_3y=x^3\) over
\(\Q(a_1,a_3)\), the discriminant, the \(3\)-division polynomial, \(-P\), the roots of \(h\), the curve
\(E'\) and the formula for \(\hat\phi\) of Section~\ref{conditions}, including \(\hat\phi\phi=[3]\).
Part~C proves the statements of Section~\ref{cyclotomic}: the Minkowski bound, class number one, the
unit group, the action of \(\zeta\mapsto\zeta^2,\zeta^3\) on \(\tau\), the norm form, the identities
for \(\pi\), the cube classes modulo \(2\); for \(j=0,1\) that \(D_j\) is smooth of genus \(1\) and
isomorphic to \(E_j\), the norm of \(\Delta(E_j)\), the splitting of \(11\), the rank bounds, the
Mordell--Weil groups with both proof flags, the map \(E_j\to D_j\to\mathbb P^1\), and that Chabauty at
\(11\) returns a bound equal to the number of points found, no unresolved residue classes, the
saturation index and the sets (9). Part~D treats Section~\ref{families}: the Tate normal form,
\(\psi_3\) and the identity for \(\psi_3(du)/d^3\); the irreducibility of (10) and its singular points;
the rank and torsion of \(W\); the birational map and its inverse; the identities (11) by an explicit
translation and shear; the cusps; the divisors of the \(t_j\); in each family the quadratic relation,
the identity \(R_j=N_j/D_j\), the integrality, degrees and palindromy of \(N_j\), the resultants, and
an explicit cube root \(g_j\) of \(N(\delta)/F_j\) (the leading coefficient is kept, as Magma's
factorization drops it); and the fibres over \(0\), \(1\), \(-1\), \(\infty\), including the residue
field, \(d\), the torsion, \(a_3/\sigma(a_3)\), \(N_{K/\Q}(\Delta)\) and \(g_j(t_0)\) at the non-cuspidal
points. Part~E carries out the checks of \hyperref[rem-both]{Remarks~2} and~\hyperref[rem-census]{3}
from the census data of 9.c, 18.c and 15.c. All computations were rerun, and the new identities of
Sections~\ref{conditions} and~\ref{families}, the exceptional fibres and the arithmetic of
Proposition~5 were cross-checked independently in Sage and PARI.

\begin{thebibliography}{P2214}

\bibitem[BL]{BL}
C. Birkenhake and H. Lange, \emph{Complex Abelian Varieties}, 2nd edition, Springer, Berlin,
2004 (Grundlehren der mathematischen Wissenschaften~302): Section~5.3, Proposition~11.4.3 and
Corollary~12.1.5.
\href{https://doi.org/10.1007/978-3-662-06307-1}{doi:10.1007/978-3-662-06307-1}.

\bibitem[Br]{Br}
N. Bruin, \emph{Chabauty methods using elliptic curves}, J. Reine Angew. Math. \textbf{562} (2003),
27--49.

\bibitem[FK]{FK}
G. Frey and E. Kani, \emph{Curves of genus 2 covering elliptic curves and an arithmetical
application}, in: Arithmetic Algebraic Geometry (Texel, 1989), Progress in Mathematics~89,
Birkh\"auser, Boston, 1991, 153--176.

\bibitem[Ka]{Ka}
S. Kamienny, \emph{Torsion points on elliptic curves and \(q\)-coefficients of modular forms}, Invent.
Math. \textbf{109} (1992), 221--229.
\href{https://doi.org/10.1007/BF01232025}{doi:10.1007/BF01232025}.

\bibitem[Ka97]{Ka97}
E. Kani, \emph{The number of curves of genus two with elliptic differentials}, J. Reine Angew. Math.
\textbf{485} (1997), 93--122. \url{https://eudml.org/doc/183547}.

\bibitem[KM]{KM}
M. A. Kenku and F. Momose, \emph{Torsion points on elliptic curves defined over quadratic
fields}, Nagoya Mathematical Journal \textbf{109} (1988), 125--149: Theorem~(0.1).
\href{https://doi.org/10.1017/S0027763000002816}{doi:10.1017/S0027763000002816}.

\bibitem[Ku]{Ku}
D. S. Kubert, \emph{Universal bounds on the torsion of elliptic curves}, Proc. London Math. Soc.
(3) \textbf{33} (1976), 193--237: Table~3.
\href{https://doi.org/10.1112/plms/s3-33.2.193}{doi:10.1112/plms/s3-33.2.193}.

\bibitem[Ma]{Ma}
B. Mazur, \emph{Modular curves and the Eisenstein ideal}, Publ. Math. IH\'ES \textbf{47} (1977),
33--186: Theorem~(8). \href{https://doi.org/10.1007/BF02684339}{doi:10.1007/BF02684339}.

\bibitem[Mi]{Mi}
J. S. Milne, \emph{On the arithmetic of abelian varieties}, Invent. Math. \textbf{17} (1972),
177--190: \S1. \href{https://doi.org/10.1007/BF01425446}{doi:10.1007/BF01425446}.

\bibitem[Na]{Na}
F. Najman, \emph{Complete classification of torsion of elliptic curves over quadratic cyclotomic
fields}, J. Number Theory \textbf{130} (2010), 1964--1968: Section~1 and Theorem~2(ii).
\href{https://doi.org/10.1016/j.jnt.2009.12.008}{doi:10.1016/j.jnt.2009.12.008};
\url{https://arxiv.org/abs/1005.0560}.

\bibitem[P30]{P30}
Census proof of 21 October 2026, \emph{An obstruction to rational \(30\)-torsion on geometrically
split abelian surfaces},
\url{https://genus-2-torsion.github.io/proofs/qsimple-30-obstruction.html}: Proposition~2,
Lemmas~3 and~4, Section~8.

\bibitem[P63]{P63}
Census proof of 21 October 2026, \emph{An obstruction to rational \(63\)-torsion on geometrically
split abelian surfaces},
\url{https://genus-2-torsion.github.io/proofs/qsimple-63-obstruction.html}: Section~1,
Lemmas~1--3 and~5--7, Sections~2--5, 7 and~9.1.

\bibitem[P2214]{P2214}
Census proof of 10 November 2026, \emph{Rational \(7\)-torsion and the \(2\)-rank of
geometrically split abelian surfaces},
\url{https://genus-2-torsion.github.io/proofs/qsimple-2214-obstruction.html}: Section~1.

\bibitem[Po]{Po}
B. Poonen, \emph{Some Diophantine equations of the form \(x^n+y^n=z^m\)}, Acta Arith. \textbf{86}
(1998), 193--205: Section~4.
\href{https://doi.org/10.4064/aa-86-3-193-205}{doi:10.4064/aa-86-3-193-205}.

\bibitem[ST]{ST}
J.-P. Serre and J. Tate, \emph{Good reduction of abelian varieties}, Ann. of Math. (2)
\textbf{88} (1968), 492--517: Theorem~1 and Corollary~2 (p.~493), Lemma~2 (p.~495), \S2
(pp.~496 and~498). \href{https://doi.org/10.2307/1970722}{doi:10.2307/1970722}.

\end{thebibliography}

\end{document}
